%%BLOCK after="Sets already detect the missing monodromy"
\begin{dcomment}{development}{p100}{The localizer of locally constant coefficients}
The pages from 97 to 126 circle around an object that the manuscript never names: the class of functors that induce an equivalence on locally constant coefficients. For a derivator it is a fundamental localizer, and it is the natural recipient of the ``2-complete'' conditions. Throughout, $\mathbb D$ is a derivator on all small categories in the sense of \citet[\S3.2]{Cisinski2009LocallyConstant} (disjoint unions go to products, evaluations at objects are jointly conservative, every restriction has both adjoints, and both Kan extensions are computed on comma categories), in the presheaf variance of the manuscript. Write $\mathrm{LC}(A)\subset\mathbb D(A)$ for the full subcategory of locally constant objects, and for $f:A\to B$ say that
\[
 f\in W^{\mathrm{lc}}_{\mathbb D}\ \text{if}\ f^*:\mathrm{LC}(B)\to\mathrm{LC}(A)\ \text{is an equivalence},\qquad
 f\in W^{1}_{\mathbb D}\ \text{if it is fully faithful};
\]
these are the 2-complete and the 1-complete equivalences of ms.~100. Let $W_{\mathbb D}$ be the class of the $\mathbb D$-equivalences (constant coefficients), and $W^{\mathrm{coh}}_{\mathbb D}$, $W^{\mathrm{hom}}_{\mathbb D}$ the classes of functors preserving cohomology, respectively homology, with every locally constant coefficient on the target. The theorem was not found in the sources consulted in this form; its proof is an application of the cited results of Cisinski, and no priority is claimed.

\textbf{Theorem.}\ours
\textit{(a) $W_\infty\subset W^{\mathrm{lc}}_{\mathbb D}\subset W^{1}_{\mathbb D}\subset W^{\mathrm{coh}}_{\mathbb D}\cap W^{\mathrm{hom}}_{\mathbb D}\subset W_{\mathbb D}$, and for a functor $A\to e$ membership in $W^1_{\mathbb D}$ and in $W_{\mathbb D}$ are the same condition.
(b) Let $A\xrightarrow{u}B\to S$ be functors. If $(u/s)^*:\mathrm{LC}(B/s)\to\mathrm{LC}(A/s)$ is an equivalence for every object $s$ of $S$, then $u^*:\mathrm{LC}(B)\to\mathrm{LC}(A)$ is an equivalence; the same holds with ``fully faithful'' in both places.
(c) $W^{\mathrm{lc}}_{\mathbb D}$ is a fundamental localizer satisfying the relative Theorem~A.
(d) $W^1_{\mathbb D}$ contains the identities, is stable under composition and retracts, satisfies the relative Theorem~A and the rule ``$f$ and $gf$ in the class imply $g$ in the class'', and in general fails two-out-of-three.}

\textit{Proof.}
(a) The inclusion $W_\infty\subset W^{\mathrm{lc}}_{\mathbb D}$ is \citet[Corollary~3.8]{Cisinski2009LocallyConstant}. If $f^*$ is fully faithful on locally constant objects then, for $M$ in $\mathbb D(e)$ and $F$ in $\mathrm{LC}(B)$, it is bijective on the morphisms from the constant object $M_B$ to $F$ and from $F$ to $M_B$; by adjunction and the Yoneda lemma in $\mathbb D(e)$ these bijections say that $f$ preserves the cohomology and the homology of $F$. Taking $F$ constant gives the last inclusion. For $A\to e$ every locally constant object of the target is constant, so full faithfulness on $\mathrm{LC}(e)=\mathbb D(e)$ is the definition of a $\mathbb D$-equivalence.

(b) \emph{Step 1: projectors.} For a subcategory $R$ of a small category $T$ let $\mathbb D(T;R)$ be the full subcategory of the objects of $\mathbb D(T)$ whose underlying presheaf inverts the arrows of $R$. The inclusion $\mathbb D(T;R)\to\mathbb D(T)$ has both adjoints. For $R=T$ this is \citet[Proposition~3.10 and Theorem~3.13]{Cisinski2009LocallyConstant}; the proof rests on Theorem~3.6 there, stated for an arbitrary set of arrows, and on Lemma~3.9, whose argument (a homotopy colimit of diagrams of spaces inverting given arrows inverts them) is the same for a subset of the arrows. Write $Q_{T,R}$ for the right adjoint.

\emph{Step 2: fibres of an opfibration.} Let $\pi:T\to S$ be a Grothendieck opfibration, $R$ the subcategory of its vertical arrows, $Q_T=Q_{T,R}$, and $i_s:T_s\to T$ the inclusion of a fibre. For $x$ over $t$ the comma category $x\backslash T_s$ is the disjoint union, over the arrows $\alpha:t\to s$, of categories with initial object the cocartesian arrow $x\to\alpha_!x$. Hence $(i_{s!}F)(x)\cong\coprod_\alpha F(\alpha_!x)$, and $i_{s!}$ sends $\mathrm{LC}(T_s)$ into $\mathbb D(T;R)$. Passing to right adjoints in the equality ``include, then $i_{s!}$'' $=$ ``$i_{s!}$, then include'' gives $i_s^*Q_T\cong Q_{T_s}i_s^*$.

\emph{Step 3: base change.} Let $v:T\to T'$ be a functor over $S$ between opfibrations which preserves cocartesian arrows. For $x'$ in $T'_s$ the inclusion $T_s/x'\to T/x'$ has a left adjoint, given by cocartesian transport. In presheaf variance restriction along a functor with a left adjoint does not change the limit, so the formula for right Kan extensions gives $i'^{*}_sv_*\cong v_{s*}i_s^*$.

\emph{Step 4: conclusion for opfibrations.} On the objects that are locally constant along the fibres there is an adjunction $v^*\dashv Q_{T'}v_*$. By Steps 2 and 3 its unit and counit restrict, on each fibre, to those of $v_s^*\dashv Q_{T'_s}v_{s*}$ on $\mathrm{LC}(T'_s)$ and $\mathrm{LC}(T_s)$. If every $v_s^*$ is an equivalence there, these are isomorphisms; every object of $T$ and of $T'$ lies in a fibre, and evaluations are jointly conservative, so $v^*$ is an equivalence between the categories of objects locally constant along the fibres. If the $v_s^*$ are only fully faithful, the same argument with the unit alone shows that $v^*$ is fully faithful. Finally $v^*$ preserves local constancy, and it reflects it when the $v_s^*$ are equivalences: for $N$ locally constant along the fibres of $T'$ and $\alpha:s\to t$, cocartesian transport gives a morphism between the restrictions of $N$ to the two fibres which becomes invertible after $v_s^*$, hence is invertible, and every arrow of $T'$ is a cocartesian arrow followed by a vertical one. So $v^*:\mathrm{LC}(T')\to\mathrm{LC}(T)$ is an equivalence, respectively fully faithful.

\emph{Step 5: reduction.} Let $\widetilde A$ be the category of triples $(a,s,\alpha:p(a)\to s)$, with $p$ the composite $A\to S$; it is an opfibration over $S$ with fibres $A/s$. The forgetful functor $\widetilde A\to A$ has the left adjoint $a\mapsto(a,p(a),\id)$, so both induce inverse equivalences on locally constant objects. The functor $u$ induces $\widetilde A\to\widetilde B$ over $S$, preserving cocartesian arrows, with fibres $u/s$. Step~4 applies.

(c) Equivalences of categories satisfy two-out-of-three and are stable under retracts, and $A\mapsto\mathrm{LC}(A)$ is a pseudofunctor; so $W^{\mathrm{lc}}_{\mathbb D}$ is weakly saturated. A category with a terminal object is $W_\infty$-aspherical, hence $W^{\mathrm{lc}}_{\mathbb D}$-aspherical by (a). The relative Theorem~A is (b).

(d) The closure properties are those of fully faithful functors, and the relative Theorem~A is the second form of (b). For the failure take sets as coefficients, $\mathbb D(A)=\operatorname{Fun}(A^{\mathrm{op}},\mathrm{Set})$, the group $G$ with two elements, and $e\xrightarrow{i}BG\xrightarrow{r}e$. The functor $r^*$ (trivial actions) is fully faithful, the composite is the identity, and $i^*$ (forgetting the action) is not full: the regular $G$-set has two equivariant endomorphisms and four endomorphisms as a set.\hfill$\square$

Three remarks. First, the compatibility of the two mate isomorphisms of Steps 2 and 3 with the units and counits, used in Step~4, is the standard one and has not been written out here. Second, Step~2 uses coproducts indexed by the arrows of $S$; the proof does not transfer as it stands to a derivator defined on finite categories only. Third, by (c) and (d) the manuscript's passage from $\mathbb D$-equivalences to 1-complete and then 2-complete ones ends at a class with all the properties of Chapter~X, whereas the intermediate classes $W^1_{\mathbb D}$, $W^{\mathrm{coh}}_{\mathbb D}$, $W^{\mathrm{hom}}_{\mathbb D}$ are not localizers. The values of these classes for sets, truncated types, spaces, complexes and spectra are tabulated in the commentary on Chapter~X, ms.~5. The axioms Der~3~a$'$, b$'$, a$''$, b$''$ of ms.~97--112 are consequences of (a) for every derivator; the axioms Der~3~c of ms.~124--125, which ask that $W_{\mathbb D}$ be contained in the smaller classes, hold for spaces and fail for each of the other coefficient theories in that table.
\end{dcomment}

%%APPEND title="Two fibrations whose composite is not a fibration"
Replacing constant-coefficient equivalences by 2-complete ones in the definition does not restore stability under composition: see the commentary at ms.~122 (``The 2-complete variant also fails to compose''). A class of fibrations that is stable under composition and base change and detects equivalences on fibres is obtained in another way, from the localizer $W^{\mathrm{lc}}_{\mathbb D}$ of the commentary at ms.~100: when that localizer is accessible, the fibrations of its Thomason model structure \citep[Th\'eor\`eme~5.2.15]{Cisinski2006} have these properties, and a functor in $W^{\mathrm{lc}}_{\mathbb D}$ between two of them over a base $S$ is in $W^{\mathrm{lc}}_{\mathbb D}$ on every fibre (both are fibrant in the slice over $S$, and the fibre functor is right Quillen there). Accessibility is known for the derivators of presentable coefficient theories only by an argument that was not checked for this commentary.

%%APPEND title="The 2-complete variant also fails to compose"
\textbf{Complements.}
(1) The same extension $E_c\to B$ shows that, with 2-complete equivalences substituted throughout, a total equivalence between fibrations over a base need not be an equivalence on the fibres: the functor $E_c\to B$ over $B$ induces an equivalence of fundamental groupoids (both total categories are simply connected), and on fibres it is $BC_2\to e$. So the assertion of ms.~84 withdrawn in footnote~41 is not repaired by the substitution. The substitution of 1-complete equivalences fails as well, by a different example which is not reproduced here. Its repair with constant-coefficient equivalences over a completely $\mathbb D$-aspherical base is not affected.
(2) The qualification ``(que je sache!)'' of ms.~122 is the reading of the edition, not a transcription slip; with the proposition above it becomes a theorem.
(3) For the homotopy fibre sequence used above one can avoid the presentation: $E_c\to B$ is a fibration and an opfibration with fibre $BC_2$ and identity transports, so $BE_c\to BB\simeq S^2$ is a fibre sequence, $\pi_1(BE_c)$ is a quotient of $C_2$, and it equals $H_1(E_c;\mathbf F_2)$, which is $0$ for the cocycle supported on one triangle and $\mathbf F_2$ for the zero cocycle.
